\section{第\ref{sec:dishbrain}章　DishBrain}

テストベンチの構成とゲームの結果は、一つの実験研究とその続編から採った。ノイズと良い設定への漂流の式は本章で導き、本書のモデルで確かめた。シャーレの中の学習の仕組みは確定していない。

{\footnotesize
\setlength{\tabcolsep}{3pt}\renewcommand{\arraystretch}{1.15}
\begin{longtable}{>{\raggedright}p{0.5cm}>{\raggedright}p{4.0cm}>{\raggedright}p{5.6cm}>{\raggedright}p{2.3cm}>{\raggedright\arraybackslash}p{1.7cm}}
\toprule
No. & 主張 & 実験と測定されたもの & 出典 & 状態 \\
\midrule\endfirsthead
\toprule
No. & 主張 & 実験と測定されたもの & 出典 & 状態 \\
\midrule\endhead
1 & テストベンチ：チップあたりマウス大脳皮質の細胞約 $800\,000$ 個またはヒト細胞約 $10^6$ 個。$8$~mm$^2$ に電極 $26\,000$ 本、記録は最大 $1024$、刺激は $8$ 本 & 論文の方法：MaxOne チップ、$8$~mm$^2$ に白金電極 $26\,000$ 本、$1024$ チャネルを毎秒 $20\,000$ サンプルで記録。技術的には最大 $32$ 電極を刺激できるが、独立に刺激できたのは $8$ 本。マウス培養はおよそ $10$ 日目から使用 & \cite{Kagan2022} & 測定済み \\
2 & $10$~ms 周期のループ：運動領域のスパイクがラケットを動かす（図\ref{fig:dishloop}） & スパイクは $10$~ms の窓ごとに数えられる。スパイクから応答刺激までの遅れは約 $5$~ms で、装置のバッファで決まる & \cite{Kagan2022} & 測定済み \\
3 & 入力：場所による符号（$8$ 本のどの電極か）と頻度 $4$--$40$~Hz & 主実験では刺激頻度が、ボールが遠い壁にあるときの $4$~Hz からラケットのそばの $40$~Hz まで線形に増える。ボールのパルスの振幅は $75$~mV、パルスは矩形の二相性 & \cite{Kagan2022} & 測定済み \\
4 & 出力 $\Delta p=c\,(n_\uparrow-n_\downarrow)$~\eqref{eq:dishreadout}、窓あたりのカウントのノイズは $1/\sqrt{RT}$ & 二つの領域の差をとる規則は論文に記述がある（活動に応じた領域の重みつき）が、正確な式はない。ノイズの推定はポアソン計数からの帰結で、本章で導いた & \cite{Kagan2022}、本章の導出 & 理論 \\
5 & 教師：ヒットの後に予測可能な刺激 $75$~mV、$100$~Hz、$0.1$~s。ミスの後に予測不能な刺激 $150$~mV、$5$~Hz、$4$~s をランダムな場所と時刻に、その後 $4$~s 休止 & 方法：ヒットの後、$75$~mV、$100$~Hz、$100$~ms を $8$ 本の電極すべてに同時に。ミスの後、$150$~mV、$5$~Hz をランダムな電極にランダムな時刻で $4$~s、その後 $4$~s 刺激なし。培養にドーパミンはない & \cite{Kagan2022} & 測定済み \\
6 & ネットワークは自分の入力が予測可能になるように行動する & 自由エネルギー原理は理論的な提案で、著者らはそこからプロトコルを導いた。実験はこれと矛盾しないが、より単純な仕組み（行~7--8）と区別できない & \cite{Friston2010,Kagan2022} & 仮説 \\
7 & 失敗がネットワークを揺さぶり成功が揺さぶらないなら、設定で過ごす時間は $\rho\propto1/P_{\text{miss}}$~\eqref{eq:dishdensity}（命題\ref{prop:dishscramble}） & 対称な揺さぶりをもつマルコフ連鎖の流れの釣り合いから従う。本章で導いた。培養での直接の検証はない & 本章の導出 & 理論 \\
8 & 「ミスでかき混ぜる」モデルは学習する：ヒットの割合 $0.21\to0.38$。毎ラリー後に揺さぶると $\approx0.12$、揺さぶりなしで $0.19$（図\ref{fig:dishmodel}） & 計算、スクリプト \texttt{models/dishbrain\_model.py}、モデル培養 $40$ 個に各 $2000$ ラリー：$0.21\to0.38$、$0.13\to0.11$、$0.19\to0.19$ & --- & モデル \\
9 & ラリーが長くなるのは完全なループの生きた培養だけ。対照は学習しない & $20$~min のセッション $399$ 回：マウス培養 $9$ 個（$101$ セッション）、ヒト $11$ 個（$138$）、培地のみのチップ $6$ 枚（$80$）、安静の培養 $20$ 個（$42$）、シミュレーション $3$ 個（$38$）。最後の $15$~min の平均ラリー長が最初の $5$~min より長いのは生きた培養だけ。ニューロンでない細胞（HEK293T）と培地のみのチップでは効果なし & \cite{Kagan2022} & 測定済み \\
10 & セッション内の有意な伸びは完全なフィードバックのときだけ。刺激の代わりに無音でも、セッションの終わりにはフィードバックなしより上手だが、効果は小さい & $3$ 日で $486$ セッション：「刺激」（$n=27$）、「無音」（$17$）、「フィードバックなし」（$15$）。時間による伸びが有意なのは「刺激」条件だけ。セッションの終わりに「無音」条件は「フィードバックなし」を小さな効果で上回った & \cite{Kagan2022} & 測定済み \\
11 & 効果は小さい。ネットワークが長期的に学習するかは未解決 & セッション内の改善は確かだが、著者ら自身によれば、数日にわたるセッション間の学習は安定しては観察されなかった & \cite{Kagan2022} & 測定済み \\
12 & プレイ中、ネットワークは臨界状態に近づき、近いほどプレイが良い & 同じ培養の活動の雪崩の解析：構造化された入力はネットワークを臨界状態に近づけ、プレイの質はそこへの近さと相関する。ただし、行動の結果についての情報がなければ、臨界性だけでは学習に足りない & \cite{Habibollahi2023} & 測定済み \\
13 & 培養の「知性」は示されていない & 30人の研究者による反論論文：閉ループでの振る舞いの変化は知性も感覚も証明しない。それを示す実験はない & \cite{Balci2023} & 仮説 \\
14 & 提案する四つ目の対照：ミスの後に同じ長さとエネルギーの予測可能な刺激（\ref{sec:dishspec}節） & この実験は行われていない。本書の提案である & --- & 仮説 \\
\bottomrule
\end{longtable}}
